% !TeX root = ../../Book3.tex
\chapter{Cauchy 定理的拓扑形式}
\label{b3:ch:10}
\BookChapterNumber{b3:ch:10}
\begin{chapterabstract}
本章完成任意平面区域中的全局 Cauchy 与留数理论。指数提取闭路的可加信息，局部原函数给同伦不变性，格网分割把零同调圈写成区域内小边界之和。得到的公式允许自交、多重绕行和多个边界分支，同时保留原函数与对数的精确周期判别。
\end{chapterabstract}

\begin{learninggoals}
\begin{itemize}
\item\label{b3:item:10-goal-1}计算闭路及整数圈对补集点的指数并使用其局部恒定性。
\item\label{b3:item:10-goal-2}用局部原函数证明同伦积分不变而不假定同伦光滑。
\item\label{b3:item:10-goal-3}证明平面零同调的格网分割与全局 Cauchy 公式。
\item\label{b3:item:10-goal-4}在一般留数公式中正确使用绕数权重并核对有限性。
\end{itemize}
\end{learninggoals}

一个围道可绕同一奇点几次，也可在不同部分反向抵消。单条简单 Jordan 边界的公式不能直接掩盖这些权重。本章把闭路升级为整数圈，用指数表达它是否在全纯区域中围住了一个合法的有限边界组合。

\ProseParagraphBreak
局部前提只有第2章的原函数与 Cauchy 理论、第3章的 Laurent 系数，以及第9章的路径同伦。格网分割是由局部走向全局的关键步骤，正文完整写出，以后不需要把“同调形式已知”藏在计算中。

% !TeX root = ../../Book3.tex
\section{指数与绕数}
\label{b3:sec:1001}

第2章用曲线收缩证明积分不变，本章再把曲线对各个孔的总绕行抽取出来。全局 Cauchy 定理最终只需要这些整数同时为零，而不要求曲线可以收缩。

\subsection{指数}
\label{b3:subsec:100101}

对分段光滑闭曲线 \(\gamma\) 及 \(a\notin\gamma\)，定义\term[quxian-zhishu]{曲线指数}[index of a curve]
\[
 \operatorname{Ind}(\gamma,a)=\frac1{2\pi\mathrm{i}}\int_\gamma\frac{\mathrm{d}z}{z-a}.
\]
沿曲线可连续提升 \(\arg(\gamma-a)\)，径向模的对数在闭路前后相同，积分只留下辐角总变化。因而指数为整数。它不是曲线某一点的局部属性，也不是闭路欧氏长度的函数。多次走同一圈使指数乘相同整数，反向使其变号。

\ProseParagraphBreak
指数作为 \(a\) 的函数在曲线外连续，整数性使它在每个补集连通分支上恒定。无界分支上它为0：当 \(|a|\) 很大时，积分由长度估计趋于0，所以这个整数最终为0。正向 Jordan 圈在内部指数1，外部0，使用的是平面分离和正向辐角变化；自交曲线则可能有多个不同整数区域。

\subsection{同伦不变性}
\label{b3:subsec:100102}

\begin{proposition}\label{b3:prop:10-index-homotopy}
在 \(\mathbb C\setminus\{a\}\) 中同伦的闭路具有相同指数。对有限整数线性组合，指数按系数相加。
\end{proposition}
\begin{proof}
可沿同伦的每条闭路连续选局部辐角，起终点差的整数在参数局部不变；紧参数区间由有限小辐角邻域覆盖，所以始终相同。若同伦路径不分段光滑，先以小盘内折线替换；圆盘内的原函数和连续辐角保证不改变指数。线性组合结论由积分线性直接得到。
\end{proof}
不经过点 \(a\) 是全部证明的条件。一个圆周连续缩成穿过 \(a\) 的点时，指数不再有定义，故不能以这样的变形宣称绕数已变为0。

\ProseParagraphBreak
指数作为补集上的局部常整数，可以在简单几何分解中直接计算。正向 Jordan 圈在内部为1、外部为0；两圈相加，在共同内部为2，在只被一圈围住处为1。若第二圈反向，共同内部为0。把两圈的公共部分画成一次边界并忽略方向，会遗漏这些权重。对自交的八字圈，两个瓣可能分别为1与 \(-1\)，取决于走向；交点本身属于曲线，指数在那里不定义。

\ProseParagraphBreak
无界补分支的指数为0也可用估计证明。曲线像落在某个有限圆盘中，当 \(|a|\) 大时，\(1/(z-a)=-a^{-1}\sum_{k\ge0}(z/a)^k\) 在曲线上一致收敛，各项全是整函数的导数，闭路积分0。再由局部常性把0传遍无界分支。这个事实保证非零指数区域不会跑到无穷远，是一般留数和只涉及有限紧填充的一个依据。

\subsection{指数与围道积分}
\label{b3:subsec:100103}

闭路积分中 \(1/(z-a)\) 专门测量绕 \(a\) 的次数。更一般有理函数的积分则是有限个这种测量的加权和，权重为留数。例如在避开0和1的区域中，\(1/(z(z-1))=-1/z+1/(z-1)\)，故
\[
 \frac1{2\pi\mathrm{i}}\int_\gamma\frac{\mathrm{d}z}{z(z-1)}
 =-\operatorname{Ind}(\gamma,0)+\operatorname{Ind}(\gamma,1).
\]
绕两个点各一次时两个贡献相消，积分为0，并不说明闭路不绕任何点。选定被积函数可能看不见某些拓扑数据；要判定全纯函数的全部周期，必须考虑整个区域及所有闭路。

\ProseParagraphBreak
指数为整数的证明可直接用连续辐角。把闭路相对于 \(a\) 归一化为 \(\eta(t)=(\gamma(t)-a)/|\gamma(t)-a|\)，有限小弧上逐段选辐角并在重叠处加 \(2\pi\) 整数以衔接，得连续 \(\theta(t)\)。闭路终点给 \(\theta(1)-\theta(0)=2\pi n\)，分段求导并积分即得 \(\operatorname{Ind}(\gamma,a)=n\)。改变最初辐角只加同一个 \(2\pi\) 整数，不改变差。这个整数解释为何随避开曲线的连续移动只能局部恒定。

\ProseParagraphBreak
对圈 \(\Gamma=\sum n_j\gamma_j\)，指数按系数相加。重走一条曲线两次给指数2，反向给负指数，完全可以出现负值。一般围道不是一个“内部”集合的正向边界；即使画出多重自交，也应逐个区域记录整数权重。Cauchy 与留数公式中的指数正是补偿这种几何复杂性的系数。

% !TeX root = ../../Book3.tex
\section{同伦形式的 Cauchy 定理}
\label{b3:sec:1002}

局部原函数使积分在路径小变形下不变。把局部相消沿同伦方格拼接，可以对一般区域的同伦路径证明积分不变；单连通性只是全部闭路都能收缩时的一个特例。

\subsection{同伦曲线}
\label{b3:subsec:100201}

\begin{theorem}\label{b3:thm:10-homotopy-integral}
设 \(f\) 在区域 \(U\) 全纯。端点固定同伦的两条分段光滑路径积分相同；自由同伦的两条闭路积分也相同。
\end{theorem}
\begin{proof}
同伦像紧，取有限个有局部原函数的圆盘覆盖它，并用 Lebesgue 数把参数方形分成足够小的格。每个小格像包含在一个圆盘内；即使内部边只是连续路径，也可用原函数端点差定义积分。每格四边端点差之和为0。把全部小格相加，内部边方向相反而抵消，只剩方形外边。端点固定情形，两条竖边常值，积分为0；闭路自由同伦情形，两条竖边是同一基点移动路径的相反方向，也抵消。因此上下两条路径积分相同。与原分段光滑路径上的普通积分一致来自原函数端点公式。
\end{proof}
这个证明明确使用全纯性产生局部原函数，不要求同伦本身可微。对连续非全纯函数没有同样结论；\(\overline z\) 在不同圆周上的积分随面积改变。

\subsection{Cauchy 积分定理}
\label{b3:subsec:100202}

\begin{corollary}\label{b3:cor:10-null-homotopic-cauchy}
若闭路 \(\gamma\) 在 \(U\) 中可收缩，\(f\) 在 \(U\) 全纯，则 \(\int_\gamma f\,\mathrm{d}z=0\)。特别地，单连通区域中所有闭路积分为零。
\end{corollary}
\begin{proof}
将 \(\gamma\) 同伦到常路径，前一定理使积分等于常路径的零积分。
\end{proof}
可收缩性对闭路本身提出强条件。区域中有孔时，某些不能收缩的闭路仍可能对全部全纯函数积分为零，例如两个孔绕行的交换子圈。Cauchy 的同调形式将精确描述积分所需的较弱条件。

\ProseParagraphBreak
同伦积分不变性适用于同一全纯函数的整个同伦像。若把围道推过一个孤立奇点，函数在某个小格上没有原函数，内部边相消仍成立，但这个格的边界积分不再为0。留数定理正是把这项缺口计为 \(2\pi\mathrm{i}\) 乘留数。于是“可自由移动围道”应读成“在避开奇点的区域里移动”，跨过奇点的变化则需另加相应指数权重。

\ProseParagraphBreak
例如 \(f(z)=1/z\)，正向小圆与正向大圆在去心平面同伦，积分都为 \(2\pi\mathrm{i}\)；但不能把其中任何一个收缩成常路径。两者的差作为圈零同调，积分为0。这一例同时展示同伦积分不变、非零周期和差圈相消，三个结论并没有相互矛盾。源路径的端点、闭路的定向及函数的定义域共同决定哪条定理可用。

\subsection{原函数判别准则}
\label{b3:subsec:100203}

\begin{theorem}\label{b3:thm:10-primitive-periods}
区域 \(U\) 上全纯 \(f\) 有单值全纯原函数，当且仅当每条分段光滑闭路的积分为0。
\end{theorem}
\begin{proof}
有原函数时由端点公式得零周期。反向固定 \(a\)，令 \(F(z)=\int_a^z f\,\mathrm{d}\zeta\)。两路径拼成闭路，零周期使定义不依赖选择。在 \(z\) 附近小盘内，增加一段线段给 \(F(z+h)-F(z)=\int_{[z,z+h]}f\,\mathrm{d}\zeta\)，除以 \(h\) 并用连续性得 \(F'=f\)。
\end{proof}
周期是对闭路同伦类的加法映射：拼接相加，反向变号，所以它落在交换群 \((\mathbb C,+)\) 中。这已经预示非交换的基本群信息会被积分压缩。若只检查一个生成圈，应先证明所有圈由它生成，例如去心平面；在多孔域没有这种简化。

\ProseParagraphBreak
无零函数 \(g\) 有全纯对数当且仅当 \(g'/g\) 的全部周期为0。若只要求全纯 \(m\) 次根，则周期可以为 \(2\pi\mathrm{i}m\) 的整数倍，因为指数除以 \(m\) 后仍为1。对数要求周期本身为零，取根只要求整除，两项条件不能混淆。

\begin{proposition}[对数与根的周期条件]\label{b3:prop:10-log-root-periods}
设 \(g\) 在区域 \(U\) 全纯且无零。它有全纯对数，当且仅当 \(\int_\gamma g'/g\,\mathrm{d}z=0\) 对全部闭路成立。它有无零全纯 \(m\) 次根，当且仅当这些周期全属于 \(2\pi\mathrm{i}m\mathbb Z\)。
\end{proposition}
\begin{proof}
若 \(g=e^L\)，则 \(g'/g=L'\)，全部周期为0。反向用零周期取原函数 \(A\)，\((g e^{-A})'=0\)，故 \(g=c e^A\)、\(c\ne0\)，取一个常数对数即可。若 \(g=h^m\)，则 \(g'/g=m h'/h\)，而 \(\int h'/h\) 是 \(h\circ\gamma\) 对0的指数乘 \(2\pi\mathrm{i}\)。反向固定 \(a\)，选 \(c^m=g(a)\)，定义
\[
 h(z)=c\exp\left(\frac1m\int_a^z\frac{g'}g\,\mathrm{d}\zeta\right).
\]
两条路径的差是闭路，假设使指数变化为 \(2\pi\mathrm{i}\) 整数，所以 \(h\) 单值；局部原函数给它全纯无零。求导得 \((h^m/g)'=0\)，在 \(a\) 值为1，故 \(h^m=g\)。
\end{proof}
例如 \(g(z)=z^m\) 在圆环上有根 \(z\)，但没有对数；其基本周期为 \(2\pi\mathrm{i}m\)。\(g(z)=z\) 在圆环上无对数也没有 \(m>1\) 次根。无零条件与区域拓扑共同决定可否选分支，不能只检查表达式中的代数形式。

\ProseParagraphBreak
在有限连通的光滑平面域中，可选各内边界的正向小绕圈作为周期测试族。用不相交割线把各内边界连到外边界，切开后成单连通域；闭路经有限多次切线交叉分解，积分等于各基本周期的整数线性组合。这项构造要求边界分支和割线有限，且确实形成切开模型。一般有无限多孔的区域不能直接宣布“检查几个孔就够了”，应保留全部闭路条件或证明所选族生成了所需的同调。

\begin{proposition}[圆环周期与 Laurent 项]\label{b3:prop:10-annulus-period}
若 \(f\) 在圆环全纯，Laurent 展开为 \(\sum_{n\in\mathbb Z}a_nz^n\)，则任意闭路满足
\[
 \int_\gamma f\,\mathrm{d}z=2\pi\mathrm{i}a_{-1}\operatorname{Ind}(\gamma,0).
\]
特别地，\(f\) 在圆环有原函数当且仅当 \(a_{-1}=0\)。
\end{proposition}
\begin{proof}
闭路像紧，落在一个较小闭圆环中，Laurent 级数在其上绝对一致收敛，故逐项积分合法。每项 \(z^n\)、\(n\ne-1\) 有单值原函数 \(z^{n+1}/(n+1)\)，贡献0；\(z^{-1}\) 项由指数定义给所写式。若 \(a_{-1}=0\)，全部周期0，由原函数判别得存在；若有原函数，正向圆圈积分0使系数为0。
\end{proof}
这条结论并非说圆环里只有一个可能的极点。函数在圆环本身没有奇点，其 Laurent 负幂描述内孔一侧的解析行为，内孔中未必有一个单独可辨的奇点。周期只看 \(a_{-1}\)，其余无穷多个负幂在积分中都不见。这又展示积分比完整解析函数保留更少的信息。

% !TeX root = ../../Book3.tex
\section{同调形式}
\label{b3:sec:1003}

同调用可加的圈代替需要顺序的路径类。在平面中，零同调可以用区域外所有点的指数为零来判断。格网分割将这个条件转换成小方形边界的相消，从而给全局积分公式。

\subsection{圈}
\label{b3:subsec:100301}

一个圈 \(\Gamma\) 是有限个分段光滑闭曲线的整数线性组合 \(\sum_j n_j\gamma_j\)。积分、指数均按系数线性定义。若先给有限条有向弧并要求每点的端点系数总和为0，可把这些弧依次拼成闭路，得到等价的圈描述；负系数可用反向弧实现。

\ProseParagraphBreak
圈允许不同部分互相抵消。相同曲线正反各走一次就是零圈；大正圈减小正圈则是圆环的正向完整边界，内边在区域边界方向中为负向。把“每条边界都正向”理解成每条都逆时针，会在多连通区域里把内边符号写错。

\ProseParagraphBreak
圈是有限的整数线性组合，而不是一条必须首尾相接的参数路径。两个彼此不相交圆周可以组成一个圈，其积分、长度上界和指数都按系数相加。把它们用一条来回连接线串成单一闭路不会改变积分和指数，因为连接线的两个方向抵消；但这只是为画图或同伦描述方便，不是圈定义的一部分。有限性则保证这些求和与曲线像的紧性可直接使用。

\ProseParagraphBreak
系数取整数是为了保留指数的拓扑含义。若只做积分恒等式，当然可以考虑复系数线性组合；然而其“指数”不再是整数权重，不能直接解释成面片覆盖次数。一般留数公式可按线性延伸处理复系数，但零同调的拓扑定义与整数圈的几何填充应先建立，再作这样的代数扩展。

\subsection{同调}
\label{b3:subsec:100302}

\begin{definition}\label{b3:def:10-planar-homology}
区域 \(U\) 中圈 \(\Gamma_1,\Gamma_2\) 称为\term[tongdiao]{同调}[homologous]，若对每个 \(a\notin U\) 有 \(\operatorname{Ind}(\Gamma_1-\Gamma_2,a)=0\)。与零圈同调的圈称为在 \(U\) 中零同调。
\end{definition}
这是平面中一维整数同调的指数判别形式。下面的格网论证也说明：零同调圈在小路径替换后确可写为有限个包含于 \(U\) 的方形边界之和。因此定义具有边界相消的几何意义，不是把积分结论先作为定义。

\ProseParagraphBreak
同伦的闭路必同调，因为区域外点不在同伦像中，指数保持；反向一般不成立。积分只依靠同调的原因是每个局部方形边界积分为零，而方形相消不保留绕各孔的先后顺序。

\ProseParagraphBreak
同调忘掉闭路的非交换次序。在双孔域中，取基于同一点的绕孔闭路 \(\alpha,\beta\)。交换子 \(\alpha*\beta*\alpha^{-1}*\beta^{-1}\) 对每个外部点的指数为0，因为四段的指数成对抵消。它因此是零同调圈，任意全纯函数积分为0。但它在双孔域中不能收缩：将区域沿割线收缩到两个圈组成的图，基本群的约化词 \(\alpha\beta\alpha^{-1}\beta^{-1}\) 不消失。图的约化词判别是这里使用的标准拓扑事实。这个例子说明同伦形式虽容易表述，仍没有覆盖一般 Cauchy 零积分的全部情况。

\ProseParagraphBreak
圈的“填充”也不能只看图中的有界区域。设 \(U=\mathbb C\setminus\{0\}\)，两个同心正向圆周的差圈 \(\Gamma=\gamma_R-\gamma_r\) 对0的指数为0，故在 \(U\) 零同调；尽管每个圆单独都围住缺失点，两者合在一起只填充中间圆环。整数面片的系数恰记录这种相消，细网格证明中的 \(n_Q\) 就是各方格内的指数。负系数是代数定向，没有对应一个“负面积的区域”。

\subsection{零同调圈}
\label{b3:subsec:100303}

\begin{lemma}[格网分割]\label{b3:lem:10-grid-cycle}
若 \(\Gamma\) 在 \(U\) 中零同调，则它可在不改变任何全纯函数积分的条件下替换成一个细格网圈，并写成有限个闭方形 \(Q\subset U\) 的整数倍边界之和。
\end{lemma}
\begin{proof}
圈的支集紧含于 \(U\)，取小于它到补集距离的格网尺度。把各路径分成位于小圆盘内的段，先以折线再以邻近格边路径替换，接点用相同短连接弧对接。所有替换都在 \(U\) 的局部圆盘内同伦，局部原函数使积分不变，区域外指数也不变。所得格圈记为 \(G\)。

对每个开格面 \(Q\)，令 \(n_Q\) 为该面上的 \(\operatorname{Ind}(G,\cdot)\)。指数在面内恒定，无界远处为0，因此非零面有限。一条有向格边上的圈系数等于其左面指数减右面指数：在边中点附近横穿一次，连续辐角的总变化正跳过该条边的有向次数；也可对小横向两点的 \(1/(z-a)\) 积分逐边直接核对。因此按格边系数逐项比较，有 \(G=\sum_Q n_Q\partial Q\)。

若某个非零面的闭包含 \(a\notin U\)，则 \(a\) 不在格圈支集上。\(a\) 的小邻域指数恒为0，因为区域外指数为0；从该邻域进入相邻格面，迫使 \(n_Q=0\)，矛盾。故每个非零闭面包含于 \(U\)，完成分割。
\end{proof}
\begin{theorem}[同调形式 Cauchy 定理]\label{b3:thm:10-cauchy-homology}
若 \(f\) 在 \(U\) 全纯，\(\Gamma\) 在 \(U\) 中零同调，则 \(\int_\Gamma f\,\mathrm{d}z=0\)。对 \(a\in U\setminus\Gamma\)，还有
\[
 \frac1{2\pi\mathrm{i}}\int_\Gamma\frac{f(z)}{z-a}\,\mathrm{d}z
 =\operatorname{Ind}(\Gamma,a)f(a).
\]
\end{theorem}
\begin{proof}
格网分割将第一积分变成有限个小方形边界积分。每个闭方形在 \(U\) 内，其凸邻域上有原函数，故积分为0。第二式把被积式写成 \(f(a)/(z-a)+(f(z)-f(a))/(z-a)\)；第二项在 \(a\) 可去且全域全纯，由第一式积分为0，第一项按指数定义给所需值。
\end{proof}

\begin{corollary}[高阶同调积分公式]\label{b3:cor:10-higher-cauchy}
设 \(\Gamma\) 在 \(U\) 零同调，\(f\) 在 \(U\) 全纯，\(a\in U\setminus|\Gamma|\)，\(n\ge0\)。则
\[
 \frac1{2\pi\mathrm{i}}\int_\Gamma
 \frac{f(z)}{(z-a)^{n+1}}\,\mathrm{d}z
 =\operatorname{Ind}(\Gamma,a)\frac{f^{(n)}(a)}{n!}.
\]
\end{corollary}
\begin{proof}
\(n=0\) 是同调 Cauchy 公式。在 \(a\) 的一个不碰曲线的圆盘内，指数为常整数。积分核及其对参数的各阶导数在紧曲线与该小盘上连续有界，所以可在积分号下微分。对原公式连续微分 \(n\) 次，核导数为 \(n!/(z-a)^{n+1}\)，遂得结论。
\end{proof}
若 \(a\) 位于指数为0的区域，全部这些积分都为0，即使 \(f(a)\) 及其导数不为0。若圈绕 \(a\) 三次，所有结果乘3。局部估计从公式取得时，要同时记录曲线长度、曲线到 \(a\) 的距离及非零指数；不能对指数为0的圈把右边除以0来求 \(f(a)\)。

\ProseParagraphBreak
Cauchy 型积分还可以在曲线外定义一个函数：\(C(w)=(2\pi\mathrm{i})^{-1}\int_\Gamma f(z)/(z-w)\,\mathrm{d}z\)。它在每个补集分支全纯，在 \(U\setminus|\Gamma|\) 等于 \(\operatorname{Ind}(\Gamma,w)f(w)\)。无界分支的指数为0，故在那里为0。跨曲线时两侧整数权重可能不同，不能因此宣称 \(C\) 在整平面全纯；其定义本来就避开曲线。这个观察将积分表示与边界跃迁区分开来。

\ProseParagraphBreak
网格填充中的两个紧性条件值得分开。曲线像必须与 \(\mathbb C\setminus U\) 有正距离，才能把曲线逐段换成细网格边而不离开定义域；非零指数面片的闭包也必须留在 \(U\)，才能对每个方格使用 Cauchy 定理。后一个条件由圈在 \(U\) 零同调及指数局部常性保证，不是前一个条件的自动后果。例如围绕去心点的圆周与缺失点距离为正，但它的填充含该点，因而不能把整张填充都当成 \(U\) 内面片。

\ProseParagraphBreak
曲线可以通过分段光滑假设出现无限多次轻微振荡，仍不能简单按每次穿网格的时刻切分，交点可能无限。本章用紧路径的有限小盘链和分段近似替换来取得有限多边形圈，替换误差圈全部落在有局部原函数的小盘内，积分与外部指数都不变。再对有限多边形应用网格分解，面片个数有限。这种顺序避免把无限网格交叉当成有限组合来求和。

\subsection{拓扑留数定理}
\label{b3:subsec:100304}

\begin{theorem}[一般留数公式]\label{b3:thm:10-residue-general}
设 \(f\) 在区域 \(U\) 亚纯，\(\Gamma\) 在 \(U\) 中零同调且不经过 \(f\) 的极点。则
\[
 \frac1{2\pi\mathrm{i}}\int_\Gamma f(z)\,\mathrm{d}z
 =\sum_{a\in U}\operatorname{Ind}(\Gamma,a)\operatorname{Res}(f,a),
\]
其中只有有限个非零项。
\end{theorem}
\begin{proof}
指数非零的有界面及其闭包组成 \(U\) 内紧集：若闭包碰到区域外点，该点附近远离圈，指数局部恒为外部的0，矛盾。离散极点集在这个紧集只有有限点 \(a_j\)。各取互不交、避开圈的小正圆 \(c_j\)，令 \(m_j=\operatorname{Ind}(\Gamma,a_j)\)。圈 \(\Gamma-\sum_jm_jc_j\) 对 \(U\) 外点及所有极点的指数都为0，故在删去极点后的区域零同调。Cauchy 定理使其积分为0，每个 \(c_j\) 上的 Laurent 系数公式给 \(\int_{c_j}f=2\pi\mathrm{i}\operatorname{Res}(f,a_j)\)，代回即得结论。
\end{proof}
同一证明允许被积函数在这些孤立点为本性奇点，只需有关非零指数点有限，并在其余区域全纯。公式对极点的限制是常见且足够的亚纯版本。

\ProseParagraphBreak
例如圆环 \(r<|z|<R\) 的完整边界在稍大圆环中零同调，内外两圈对中心的指数相消；对全纯函数积分为0。但单独外圈不零同调，\(1/z\) 的积分非零。不能因圈“看起来是边界”便省去检查它是否是当前全纯区域中的边界。

\begin{corollary}[一般围道的辐角原理]\label{b3:cor:10-argument-general}
设 \(f\) 在 \(U\) 亚纯且不恒为0，\(\Gamma\) 在 \(U\) 零同调并避开 \(f\) 的零点与极点。则
\[
 \operatorname{Ind}(f\circ\Gamma,0)=
 \sum_a\operatorname{Ind}(\Gamma,a)\operatorname{ord}_a f.
\]
左边对圈逐段按整数系数理解，右边只有有限个非零项。
\end{corollary}
\begin{proof}
\(f'/f\) 在零极点处的留数为 \(\operatorname{ord}_a f\)。一般留数定理给右侧等于 \((2\pi\mathrm{i})^{-1}\int_\Gamma f'/f\,\mathrm{d}z\)；分段链式法则把它变成像圈的指数。零极点在相关的非零指数紧填充中只有有限多个，所以和有限。
\end{proof}
如果圈含负系数，式子是带权零极点计数。比如两个嵌套圆的差只计数圆环中的零极点，内孔的权重为0。这比只说“圈内的零点数减极点数”精确，因为一般圈没有单一的内部。

\ProseParagraphBreak
还可以在核前乘一个全纯函数 \(h\)，得到加权零极点和
\[
 \frac1{2\pi\mathrm{i}}\int_\Gamma h(z)\cdot\frac{f'(z)}{f(z)}\,\mathrm{d}z
 =\sum_a\operatorname{Ind}(\Gamma,a)\cdot h(a)\cdot\operatorname{ord}_a f.
\]
选 \(h(z)=z^k\) 就计算带重数零极点的幂和。若只有一个指数为1的 Jordan 围道且里面无极点，\(k=1\) 给零点总和，\(k=0\) 给零点总数。计算需避开围道零点，不是对含零点的边界积分取未声明的主值。

\ProseParagraphBreak
割线围道中的两岸方向尤其容易出错。对正实轴割开的平面取 \(0<\arg z<2\pi\)，上岸由内半径走向外半径，下岸由外走向内；两岸的 \(z^{s-1}\) 因子相差 \(e^{2\pi\mathrm{i}(s-1)}\)。因此围道积分中两岸合成为 \((1-e^{2\pi\mathrm{i}s})\int_r^R x^{s-1}/(1+x)\,\mathrm{d}x\)，而不是两项直接相加。小圆弧的估计为 \(O(r^{\operatorname{Re}s})\)，大圆弧为 \(O(R^{\operatorname{Re}s-1})\)，恰需 \(0<\operatorname{Re}s<1\)。位于 \(-1\) 的留数按分支取 \(e^{\pi\mathrm{i}(s-1)}\)。整理得
\[
 \int_0^\infty\frac{x^{s-1}}{1+x}\,\mathrm{d}x=\frac\pi{\sin\pi s}.
\]
这个计算先在两岸略离割线的有限围道使用留数定理，再令岸逼近、半径趋于0和无穷；各项估计说明极限合法。后面特殊函数中会再次出现此积分，届时使用的是已确认的分支与收敛范围。

\ProseParagraphBreak
一般留数定理并不要求整个 \(U\) 只有有限个极点。相关圈的非零指数区域被一个紧子集包住，亚纯函数的极点在紧子集中只有有限个；圈外的其他极点权重为0，不出现在和中。若极点向 \(U\) 的边界聚集，仍与局部亚纯性相容。若它们在 \(U\) 内聚集，就不能是通常意义的亚纯函数，公式的假设已失败。

\ProseParagraphBreak
留数在小圈上的计算与大圈上的拓扑权重分工明确。局部 Laurent 系数决定每个奇点贡献多少，指数决定围道取该贡献几次。先局部计算再乘权重，可处理嵌套圈、负向圈和多重绕圈；如果从图像直接判断一个奇点“在里面”，容易漏掉权重0、负权重或大于1的情形。计算过程还应保留奇点恰落在曲线上的例外，此时通常积分未定义，需要改变围道或另行声明主值问题。

\ProseParagraphBreak
拓扑条件也不能用“所有奇点都有指数0”在任意函数上替代。对一个特定亚纯函数，某些缺失点未必是它的奇点；即使其已知奇点的加权留数为0，函数也可能在未填入的定义域外没有任何指定延拓。一般 Cauchy 定理的零同调要求针对整个 \(\mathbb C\setminus U\)，保证填充合法；针对某一函数的特殊抵消可以给额外零积分，但应通过该函数的实际延拓或直接计算证明。

\ProseParagraphBreak
本章从路径同伦走向整数圈同调，没有引入高维同调的全部定义。这里的指数判别与平面面片填充已经足够证明需要的 Cauchy 公式。一般流形上的链、边界算子与同调群会把这个过程抽象化；若在曲面上没有全局平面坐标，就不能继续以“外部点的绕数”直接定义同调，必须改用那套链理论。附录A的曲面留数证明通过局部三角剖分相消，正是保持同一积分思想并换成适合曲面的对象。


\begin{chaptersummary}
同伦保留路径顺序，同调保留区域外各点的指数。全纯积分是可加量，故零同调已经足够使其消失。一般 Cauchy 公式与留数公式均带指数权重；有限性来自非零指数区域紧含于全纯或亚纯区域。
\end{chaptersummary}

\BookExercises

% 本章题组自拟；以正文已建立的工具设计计算、反例、构造与综合训练。
\begin{exercise}\label{b3:ex:10-1}
外正圆半径3减去内正圆半径1构成圈。求它对0、2、4的指数。
\end{exercise}
\begin{solution}
对0是 \(1-1=0\)，对2是 \(1-0=1\)，对4是0。这个圈是圆环完整正向边界，内边的负系数不可漏掉。
\end{solution}

\begin{exercise}\label{b3:ex:10-2}
计算同一正圆重复三次再反向一次，对圆内点的指数。
\end{exercise}
\begin{solution}
线性性给 \(3-1=2\)。对应的 \(1/(z-a)\) 积分是 \(4\pi\mathrm{i}\)，长度本身不决定该整数。
\end{solution}

\begin{exercise}\label{b3:ex:10-3}
设 \(f\) 在 \(|z|<4\) 全纯，上题的圈为 \(\Gamma\)。计算 \(\int_\Gamma f(z)/(z-2)\,\mathrm{d}z\)。
\end{exercise}
\begin{solution}
这里使用第一题的圆环圈。它在半径4圆盘零同调，对2指数1，故积分为 \(2\pi\mathrm{i}f(2)\)。若分母改为 \(z\)，指数0，积分为0。
\end{solution}

\begin{exercise}\label{b3:ex:10-4}
证明闭路拼接的交换子 \(\alpha*\beta*\alpha^{-1}*\beta^{-1}\) 对任意外点指数为0，且全纯积分为0。
\end{exercise}
\begin{solution}
指数与积分均按拼接相加，反向变号，四项分别抵消。因此该圈在区域中零同调，并有零积分；这不以交换子是否可收缩为前提。
\end{solution}

\begin{exercise}\label{b3:ex:10-5}
在去掉0和1的平面中，圈对两点指数为2与 \(-1\)。求 \(\int_\Gamma 1/(z(z-1))\,\mathrm{d}z\)。
\end{exercise}
\begin{solution}
部分分式系数为 \(-1,1\)，故积分 \(2\pi\mathrm{i}(-2-1)=-6\pi\mathrm{i}\)。此计算直接由两项指数定义，也可把圈放在包含0和1的亚纯区域用一般留数公式。
\end{solution}

\begin{exercise}\label{b3:ex:10-6}
证明 \(f\) 在去心平面全纯时有原函数当且仅当单位圆积分为0。
\end{exercise}
\begin{solution}
每条闭路的同伦类由整数绕数决定，第9章已完整计算基本群。积分同伦不变且拼接相加，所以任意周期是单位圆周期的整数倍；该周期为0时全部周期为0，再用原函数判别。
\end{solution}

\begin{exercise}\label{b3:ex:10-7}
设 \(g\) 在区域无零。说明 \(g\) 有对数与有平方根各自对 \(g'/g\) 周期的要求。
\end{exercise}
\begin{solution}
对数要求每个周期为0。平方根要求每个周期属于 \(4\pi\mathrm{i}\mathbb Z\)，使 \(\exp(\frac12\int g'/g)\) 路径独立；初值取 \(g(a)\) 的任一平方根即可。必要性由 \(g=h^2\) 的对数导数和绕数得到。
\end{solution}

\begin{exercise}\label{b3:ex:10-8}
若圈在 \(U\) 零同调，是否在较小包含其支集的 \(V\subset U\) 中仍零同调？给反例。
\end{exercise}
\begin{solution}
不一定。单位正圆在全平面零同调，却在去心平面不零同调，因为新补集点0的指数为1。删去一个点可能使原来的填充方形不再属于区域，因此全纯积分结论的区域不能任意缩小。
\end{solution}
